% No preamble. One-step quantum interface. Prefix os-.
\section{One fixed-ratio fluctuation step in physical coordinates}
\label{os-section}

We work in dimension four with a fixed odd blocking ratio
$L\ge3$. Fine spacing is $\eta>0$ and the next retained
spacing is $h=L\eta$. The number of retained cells may be
arbitrarily large; constants may depend on $L$ and the gauge
group, but not on $h$, $\eta$, or compatible periodic volume.
This is one blocking step, rather than a simultaneous
integration of all fluctuations below a fixed macroscopic scale.

Let $y=x/h$, and write the rescaled potential as
$A_s(y)=hA_{\rm phys}(x)$. Then
\begin{equation}\label{os-rescaling}
 U(e)=e^{i\eta A_{\rm phys}(e)}=e^{iA_s(e)/L},\qquad
 \nabla_\eta A_{\rm phys}=h^{-2}\nabla_{1/L}A_s.
\end{equation}
The four-dimensional coupling-free Wilson action is unchanged
by this relabeling: it is the same sum of plaquette densities,
each counted once. The physical Maxwell action likewise obeys
\begin{equation}\label{os-quadratic-scaling}
 \eta^4\sum|d_\eta A_{\rm phys}|^2
       =L^{-4}\sum|d_{1/L}A_s|^2.
\end{equation}
The corresponding scalar gauge projection and flat averaging
constraint are the fixed-$L$ unit-coarse-scale operators of
the preceding notebooks after this rescaling.

Write $X_L=\ker Q_L\cap\ker(R_LD_*^T)$, where
$D_*=D_{1/L}$, and let $\mathcal C_L$ be the Maxwell
variational inverse on $X_L$. For the physical constrained
Gaussian $Z_{\rm phys}$ and its rescaling $Z_s=hZ_{\rm phys}$,
the covariance \emph{operators} under the physical pairings
and their raw coordinate matrices satisfy
\begin{equation}\label{os-covariance-scaling}
 \mathcal C_h=h^2\mathcal C_L,\qquad
 \operatorname{Cov}_{\rm count}(Z_{\rm phys})
       =\eta^{-4}h^2\mathcal C_L,\qquad
 \operatorname{Cov}_{\rm count}(Z_s)=L^4\mathcal C_L.
\end{equation}
The last two are coordinate covariances, rather than operators
under weighted physical sums. They follow also by pushing the
normalized Gaussian through \eqref{os-quadratic-scaling};
the constant Lebesgue Jacobian cancels from its normalization.

Define the complete dimensionless value/first-difference vector
\begin{equation}\label{os-microscopic-field}
 Y=\bigl(\eta Z_{\rm phys},\eta^2\nabla_\eta Z_{\rm phys}\bigr)
   =\bigl(Z_s/L,\nabla_*Z_s/L^2\bigr).
\end{equation}
Equation \eqref{os-covariance-scaling} gives the exact counting
covariance
\begin{equation}\label{os-Y-covariance}
 \Sigma_Y=S_L\mathcal C_LS_L^T,
       \qquad S_L=(LI,\nabla_*).
\end{equation}
The previously proved energy sandwiches imply
\begin{equation}\label{os-uniform-covariance}
 0\le\Sigma_Y\le K_L I
\end{equation}
for a fixed finite $K_L$, uniformly in spacing and volume.
Indeed $\mathcal C_L\le G_{0,L}$ and
$(I,\nabla_*)G_{0,L}(I,\nabla_*)^T$ is uniformly bounded;
the extra factor $L$ is fixed. All real trace-orthonormal
color components can be included at once.

Every retained cell and a fixed collar contain a fixed number
$m_L$ of these scalar tests. A test belongs to at most a fixed
number $\nu_L$ of such collars. Consequently a covariance-based
joint-tail theorem applies with constants independent of the
physical scale and the number of cells. Singularity of
$\Sigma_Y$ from gradients or constraints causes no problem for
such an upper bound.

\subsection{Small nonconstant backgrounds: the actual reduced Hessian}
\label{os-nonflat}

We can also control the quadratic fluctuation form at the
nonlinear corrected minimizer from Notebook 15, on its stated
small global coarse-data ball. Everything in this subsection
is first expressed at retained spacing one, with fine spacing
$1/L$. Let
\[
 A=\xi+H[B-\mathcal R(A)],\qquad
 A_*=A(\xi_*(B),B),\qquad
 T_A=(I+HD\mathcal R(A))^{-1}.
\]
Use the physical trace pairing and let $K_X$ be the compression
of $K=d^Td$ to $X_L$. Write $\mathcal B_B$ for the Hessian,
in the linear $\xi$ coordinate, of the exact reduced Wilson
action at $\xi_*(B)$. This is a fine fluctuation Hessian;
it is different from the coarse Hessian $D_B^2V$ of Notebook 16.

For $z,w\in X_L$, put $v=T_{A_*}z$ and $u=T_{A_*}w$.
The exact second derivative of the constraint chart gives
\begin{equation}\label{os-reduced-form}
 \mathcal B_B[z,w]
  =\langle v,K u\rangle
   +\langle S_{\rm jet}v,W(A_*)S_{\rm jet}u\rangle
   -\langle m(B),D^2\mathcal F(A_*)[v,u]\rangle.
\end{equation}
Here $W(A_*)=Df_{1/L}(S_{\rm jet}A_*)$ has norm
$O(\norm{B}_\infty)$ and $m(B)=O(\norm{B}_\infty)$.
The last term is retained; a Hessian of the ambient action
alone would not be the reduced fluctuation form.

The previously established tangent estimate is
\begin{equation}\label{os-tangent-energy}
 \norm v_2+\norm{\nabla_*v}_2
       \le C\langle z,Kz\rangle_{\rm Herm}^{1/2},
\end{equation}
and the same holds for $u,w$, also on the common complex
background ball. To see its complex extension directly, use
$v=z-HD\mathcal R(A_*)v$, the physical Hilbert smallness
$\norm{D\mathcal R(A_*)}_{2\to2}=O(\norm B_\infty)$,
the bounded $H,\nabla_*H$, and the flat coercive bound on $X_L$.

The flat mixed terms vanish, giving the exact bilinear identity
\begin{equation}\label{os-K-difference}
 \langle v,Ku\rangle-\langle z,Kw\rangle
   =\langle D\mathcal R(A_*)v,
                      q D\mathcal R(A_*)u\rangle.
\end{equation}
The right side is $O(\norm B_\infty^2)$ relative to the two
Maxwell energy norms. The middle term of
\eqref{os-reduced-form} is $O(\norm B_\infty)$ by
\eqref{os-tangent-energy}. The genuine two-Hilbert-direction
constraint estimate bounds the last term by the same order.
Consequently, on a fixed smaller complex ball,
\begin{equation}\label{os-relative-bound}
 \mathcal B_B=K_X^{1/2}(I+E_B)K_X^{1/2},\qquad
 E_B^T=E_B,\qquad \norm{E_B}_{2\to2}\le C\norm B_\infty.
\end{equation}
The normalized matrix $E_B$ is holomorphic in $B$. The square
root is that of the fixed real positive $K_X$; the norm is
the Hermitian operator norm, while the displayed transpose
is the analytic bilinear transpose.

For real $B$, choose the radius so that $\norm{E_B}\le\epsilon<1$.
The centered Gaussian $\zeta_s\in X_L$ for this reduced
quadratic form has covariance operator
$\mathcal C_B=J_X\mathcal B_B^{-1}J_X^T$ and satisfies
\begin{equation}\label{os-real-comparison}
 (1+\epsilon)^{-1}\mathcal C_L
     \le\mathcal C_B\le(1-\epsilon)^{-1}\mathcal C_L.
\end{equation}
Its dimensionless \emph{chart-coordinate} vector
$(\zeta_s/L,\nabla_*\zeta_s/L^2)$ therefore has counting
covariance bounded by $(1-\epsilon)^{-1}K_LI$.

The linearized physical-potential fluctuation is
$T_{A_*}\zeta_s$. It also has a uniformly bounded
dimensionless value/gradient covariance. Indeed
\eqref{os-tangent-energy} bounds
$S_LT_{A_*}J_XK_X^{-1/2}$ uniformly, so
\begin{equation}\label{os-tangent-covariance}
 \norm{S_LT_{A_*}\mathcal C_BT_{A_*}^TS_L^T}_{2\to2}
       \le C_L(1-\epsilon)^{-1}.
\end{equation}
Here the finite-dimensional maps on $X_L$ are read with the
same inclusions as in \eqref{os-real-comparison}. This is a
real Gaussian statement at a real background.

For complex $B$, \eqref{os-relative-bound} instead defines a
normalized complex Gaussian functional with strictly positive
Hermitian real part of its precision. Estimates for fixed real
coordinate tests can be obtained from its finite-dimensional
marginals. A cutoff whose defining tests themselves depend on
complex $B$ is a different problem; no holomorphy is inferred
by placing an indicator around such moving complex tests.

\subsection{The local domain for the nonlinear chart}
\label{os-chart-domain}

At the fixed mesh $1/L$, bounded values control all physical
first differences and their H\"older seminorms with constants
depending only on $L$ and $\alpha$. Thus on the event that
\emph{all fine value coordinates in the regulator} of the
dimensionless chart vector have modulus at most $R$,
\begin{equation}\label{os-domain-radius}
 \norm{g\zeta_s}_{N_\alpha}\le C_{L,\alpha}gR.
\end{equation}
Choose $gR$ small enough to stay inside the nonlinear chart
and its analytic action domain. If $R(g)=(\log(1/g))^p$
for fixed $p>0$, then both $gR(g)$ and $gR(g)^3$ tend to
zero. This is compatible with a large Gaussian-tail threshold.
The same norm comparison on a fixed cell and its stencil
collar is a local comparison of fields. It does not extend
the globally defined nonlinear minimizer chart through an
arbitrary complement of bad cells; that requires a regional
construction with its own boundary estimates.
It is the relevant fixed-ratio local scaling, rather than
a uniform smooth bound on all unresolved macroscopic modes.

On arbitrarily large volumes the event that \emph{every} cell
is good need not have large probability. The quantum step
requires joint bounds on prescribed bad cells and a summation
over their connected regions. The covariance bound supplies
the first of these ingredients. Integrating the nonlinear
remainder, the Haar and coordinate Jacobians, and bad regions
remains a further task; the present Gaussian comparison does
not replace those integrals.
